% Section 2. Budget 2.5 pages.
% Source: draft/bie-enclosure/s-modelproblem.tex (clean version).
\section{Model problem and the common gap}\label{sec:model}

\subsection{Periodic media with a line defect}
For $s\in\{-,+\}$ let $\rho_s$ be a coefficient with periods $L_s>0$ in $x$ and
$d>0$ in $y$; the two sides share the period $d$, while $L_-$ and $L_+$ may
differ. Each $\rho_s$ is piecewise constant: in a common background
$\rho_{\mathrm e}>0$, every periodic cell contains exactly one simply connected
rod with $C^2$ boundary and constant value $\rho_{s,\mathrm c}>0$. The
coefficient $\rho$ of the waveguide equals $\rho_-$ for $x<x_-$, $\rho_+$ for
$x>x_+$, and in the center cell $(x_-,x_+)\times I_d$, $x_\pm=\pm L_0/2$,
contains one further rod $\Omega_0$ or none. Here $I_d:=(-d/2,d/2)$, and the
cell windows are chosen so that every rod closure stays away from the cell
boundary. We write $\Upsilon$ for the union of all rod boundaries,
$0<\rin\le\rho\le\rax$, and $B:=\RR\times I_d$.

Fix $\beta\in[-\pi/d,\pi/d)$. A guided mode at wavenumber $k>0$ is a nonzero
solution of $-\Delta w=k^2\rho w$ in $\RR^2$ of the form
$w=v\,\ee^{\ii\beta y}$ with $v$ $d$-periodic in $y$ and $v|_B\in H^1(B)$
\cite{Fliss2013}; on $\Upsilon$ the value and the normal derivative of $w$ are
continuous. Let $H^1_\beta(B)$ be the functions in $H^1(B)$ whose traces on the
upper and lower edges $\Sigma^\uparrow,\Sigma^\downarrow$ satisfy
$u|_{\Sigma^\uparrow}=\ee^{\ii\beta d}u|_{\Sigma^\downarrow}$. On
$\mathcal H:=L^2(B,\rho\,\dd x\,\dd y)$ define, following
\cite[sect.~3]{Fliss2013},
\begin{equation}\label{eq:operator}
  A(\beta)u:=-\rho^{-1}\Delta u,\quad
  D(A(\beta)):=\bigl\{u\in H^1_\beta(B):\Delta u\in L^2(B),\
  \partial_y u|_{\Sigma^\uparrow}=\ee^{\ii\beta d}\partial_yu|_{\Sigma^\downarrow}\bigr\}.
\end{equation}
The condition $\Delta u\in L^2(B)$ excludes distributional sources on $\Upsilon$
and thus contains the transmission conditions. $A(\beta)$ is the nonnegative
self-adjoint operator of the closed form
$a[u,v]=\int_B\nabla u\cdot\overline{\nabla v}$, $D(a)=H^1_\beta(B)$, in
$\mathcal H$, and a guided mode is an eigenpair $A(\beta)u=k^2u$. Replacing
$\rho$ by $\rho_s$ in the whole strip gives the periodic operators $A_s(\beta)$
on $\mathcal H_s:=L^2(B,\rho_s\,\dd x\,\dd y)$. Since $\rho$ enters only the
mass term, Fourier analysis on the quasi-periodic cylinder gives
\begin{equation}\label{eq:domain-h2}
  D(A(\beta))=D(A_\pm(\beta))=H^2_\beta(B),
\end{equation}
with equivalent graph and $H^2$ norms. In particular, functions in
$D(A(\beta))$ have no jumps across $\Upsilon$; the trace estimate of
\cref{sec:bound} rests on this fact.

\subsection{Essential spectrum and the common gap}
For $\theta\in[-\pi/L_s,\pi/L_s)$ let $A_s(\beta,\theta)$ be the Bloch operator
on the cell $(0,L_s)\times I_d$ with phases $\ee^{\ii\theta L_s}$ in $x$ and
$\ee^{\ii\beta d}$ in $y$, and let $\lambda_{s,1}(\theta)\le\lambda_{s,2}(\theta)
\le\cdots$ be its eigenvalues. Here $\theta$ is the Bloch wavenumber of the lead
along $x$, not a second transverse phase: a band that crosses an energy at some
$\theta$ carries a wave into that lead.

\begin{proposition}[essential spectrum with two leads]\label{prop:ess}
For every real $\beta$,
\[
  \sigma_{\mathrm{ess}}(A(\beta))=\sigma(A_-(\beta))\cup\sigma(A_+(\beta)),
  \qquad
  \sigma(A_s(\beta))=\bigcup_{n\ge1}\bigcup_{\theta}\{\lambda_{s,n}(\theta)\}.
\]
\end{proposition}

For identical leads this is \cite[Prop.~3.1]{Fliss2013}; \cite[Rem.~2.1(2)]{Fliss2013}
notes the extension to different leads. A proof (a parametrix built from the two
periodic resolvents, and Weyl sequences supported far out in either lead) is
given in \cref{sm:proof-ess}. Since $\|\nabla u\|^2\ge\|\partial_y u\|^2\ge
b_\beta^2\|u\|^2_{L^2(B)}$ with $b_\beta:=\min_m|\beta+2\pi m/d|$, all three
operators are bounded below by $\lambda_0:=b_\beta^2/\rax$, which is positive
for $\beta\notin2\pi\ZZ/d$.

\begin{assumption}[common gap]\label{ass:gap}
For fixed $\beta$, the open interval $G_\beta=(g_-,g_+)$, $0<g_-<g_+$, does not
meet $\sigma(A_-(\beta))\cup\sigma(A_+(\beta))$.
\end{assumption}

\begin{proposition}[discrete spectrum in the gap]\label{prop:discrete}
Under \cref{ass:gap} the set $\sigma(A(\beta))\cap G_\beta$ consists of isolated
eigenvalues of finite multiplicity, which can accumulate only at $g_\pm$; the
eigenfunctions decay exponentially as $|x|\to\infty$.
\end{proposition}

The proof (\cref{sm:proof-discrete}) uses a conjugated resolvent of $A_\pm$
with a weight that vanishes on the support of the source. The spectral-distance
bound of \cref{sec:bound} holds for all of $\sigma(A(\beta))$; the gap is used
only to interpret the spectral point it produces.

\subsection{A finite-dimensional test for the common gap}\label{sec:gap-test}
To conclude that an interval $[\mu-R,\mu+R]$ contains an eigenvalue, we must
show for each lead that the $n_s$th band lies below the interval and the
$(n_s+1)$st band above it, for all $\theta$. Two ingredients reduce this to
finitely many matrix tests.

\begin{lemma}[Lipschitz bound for band roots]\label{lem:lipschitz}
If $\rho_s\ge\rho_{\min,s}>0$, then for all $n$, $\theta$, $\theta'$,
$|\lambda_{s,n}(\theta)^{1/2}-\lambda_{s,n}(\theta')^{1/2}|
\le|\theta-\theta'|/\rho_{\min,s}^{1/2}$.
\end{lemma}

\begin{proof}
After the gauge change $v=\ee^{-\ii\theta x}u$ the fiber form is
$a_\theta[v]=\|(\partial_x+\ii\theta)v\|^2+\|\partial_yv\|^2$ on a form domain
independent of $\theta$, and $|a_\theta[v]^{1/2}-a_{\theta'}[v]^{1/2}|\le
|\theta-\theta'|\,\|v\|_{L^2}\le|\theta-\theta'|\,\|v\|_{\mathcal H_s}/
\rho_{\min,s}^{1/2}$. Apply the min--max principle to the square root of the
Rayleigh quotient.
\end{proof}

The upper band needs a guaranteed \emph{lower} bound, which Ritz values do not
give. On the gauged cell take the orthonormal plane waves
$\psi_{p,q}=\ee^{2\pi\ii px/L_s}\ee^{\ii\beta_qy}/\sqrt{L_sd}$,
$\beta_q=\beta+2\pi q/d$, with symbols
$d_{p,q}(\theta)=(2\pi p/L_s+\theta)^2+\beta_q^2$. For a finite index set
$\mathcal I$ let $\Lambda_{\mathcal I}(\theta)=\inf_{(p,q)\notin\mathcal I}
d_{p,q}(\theta)$, $K_{\mathcal I}=\diag_{\mathcal I}d_{p,q}(\theta)$,
$M_{\mathcal I}=(\langle\rho_s\psi_j,\psi_i\rangle)_{i,j\in\mathcal I}$, and
$G_{\mathcal I}=(\langle\rho_s^2\psi_j,\psi_i\rangle)-M_{\mathcal I}^2$, the Gram
matrix of the components of $\rho_s\psi_j$ outside the span.

\begin{theorem}[Fourier complement comparison]\label{thm:complement}
Let $\lambda\ge0$ and $c_{\mathcal I}:=\Lambda_{\mathcal I}(\theta)-
\lambda\rho_{\max,s}>0$. With $T_{\mathcal I}:=K_{\mathcal I}-\lambda M_{\mathcal I}$
and $T^-_{\mathcal I}:=T_{\mathcal I}-\lambda^2G_{\mathcal I}/c_{\mathcal I}$, and
$n(\cdot)$ the number of negative eigenvalues,
\[
  n(T_{\mathcal I})\le\#\{j:\lambda_{s,j}(\theta)<\lambda\}\le n(T^-_{\mathcal I}).
\]
\end{theorem}

The left inequality is the Ritz bound; the right one follows from coercivity on
the complement after completing the square in the cross term
(\cref{sm:proof-complement}). Keeping the matrix $G_{\mathcal I}$, rather than a
scalar contrast bound, is what makes a lower bound possible with a small basis
($|\mathcal I|=169$ below). For a rod $\Omega_s$ with contrast
$\delta_s=\rho_{s,\mathrm c}-\rho_{\mathrm e}$ and
$B_{\mathcal I}=(\langle1_{\Omega_s}\psi_j,\psi_i\rangle)$ one has
$M_{\mathcal I}=\rho_{\mathrm e}I+\delta_sB_{\mathcal I}$ and
$G_{\mathcal I}=\delta_s^2(B_{\mathcal I}-B_{\mathcal I}^2)$; the entries of
$B_{\mathcal I}$ are Fourier coefficients of the rod, which reduce to periodic
boundary integrals.

\begin{proposition}[panel test]\label{prop:panel}
Cover $[-\pi/L_s,\pi/L_s]$ by panels with centers $\theta_i$ and radii $r_i$,
and let $0<k_L<k_U$, $\tilde r_i=r_i/\rho_{\min,s}^{1/2}<k_L$. If at every center
\begin{equation}\label{eq:panel-tests}
  n\bigl(T_{\mathcal I}((k_L-\tilde r_i)^2,\theta_i)\bigr)\ge n_s,\qquad
  n\bigl(T^-_{\mathcal I}((k_U+\tilde r_i)^2,\theta_i)\bigr)\le n_s,
\end{equation}
with $c_{\mathcal I}>0$ in the second test, then $\lambda_{s,n_s}(\theta)<k_L^2$
and $\lambda_{s,n_s+1}(\theta)\ge k_U^2$ for all $\theta$.
\end{proposition}

\begin{proof}
\Cref{thm:complement} gives $\lambda_{s,n_s}(\theta_i)<(k_L-\tilde r_i)^2$ and
$\lambda_{s,n_s+1}(\theta_i)\ge(k_U+\tilde r_i)^2$; by \cref{lem:lipschitz} the
band roots move by at most $\tilde r_i$ within a panel.
\end{proof}

The test also certifies the band index $n_s$. If the rod is symmetric under
$y\mapsto-y$, then $\lambda_{s,n}(\theta)=\lambda_{s,n}(-\theta)$ and the panels
need to cover only $[0,\pi/L_s]$. Writing $\overline a_s=k_L^2$ and
$\underline b_s=k_U^2$ for the certified values of lead $s$, the interval
$[\mu-R,\mu+R]$ lies in a common gap if
\begin{equation}\label{eq:gap-condition}
  \max(\overline a_-,\overline a_+)<\mu-R
  \quad\text{and}\quad
  \mu+R<\min(\underline b_-,\underline b_+).
\end{equation}
A rounding-inclusive implementation must also bound the quadrature error in
$B_{\mathcal I}$ (\cref{sm:band-tests}).
